\begin{table}[htbp]
\centering
\footnotesize
\caption{Linearidade cruzada (google/gemini-2.5-flash-lite): fidelidade e acurácia em R$^2$, R$^3$ e R$^4$ nos problemas linear e não lineares.}
\label{tab:final_linearidade_cruzada_R2_R3_R4_flashlite}
\begin{tabular}{lllrrrrrrrr}
\toprule
 &  &  & \multicolumn{4}{c}{Fidelidade (vs LLM) (\%)} & \multicolumn{4}{c}{Acur.\ vs real (\%)} \\ \cmidrule(lr){4-7} \cmidrule(lr){8-11}
Problema & Atributos & Algoritmo & R$^2$ & R$^3$ & R$^4$ & $\Delta$ R$^4$$-$R$^2$ & R$^2$ & R$^3$ & R$^4$ & $\Delta$ \\
\midrule
A\_linear & — & perceptron & 81,3 & 93,1 & 83,8 & +2,4 & — & — & — & — \\
A\_linear & — & nnls & 82,7 & 92,2 & 84,2 & +1,6 & — & — & — & — \\
\midrule
F\_meia\_lua & — & perceptron & 81,6 & 85,1 & 78,1 & $-$3,5 & 63,5 & 61,9 & 63,2 & $-$0,3 \\
F\_meia\_lua & — & nnls & 81,3 & 84,1 & 81,3 & +0,0 & 65,1 & 64,8 & 65,1 & +0,0 \\
\midrule
G1\_homem\_mulher & — & perceptron & 76,9 & 78,1 & 78,6 & +1,7 & 42,6 & 41,4 & 41,4 & $-$1,2 \\
G1\_homem\_mulher & — & nnls & 78,6 & 82,9 & 84,5 & +6,0 & 37,1 & 36,7 & 43,1 & +6,0 \\
\midrule
G1\_homem\_mulher\_classesAB & — & perceptron & 91,9 & 91,7 & 90,5 & $-$1,4 & 13,1 & 13,3 & 14,0 & +1,0 \\
G1\_homem\_mulher\_classesAB & — & nnls & 94,3 & 93,6 & 93,6 & $-$0,7 & 16,0 & 16,2 & 15,7 & $-$0,2 \\
\midrule
G2\_homem\_mulher & — & perceptron & 87,1 & 87,6 & 89,5 & +2,4 & 78,3 & 78,3 & 76,4 & $-$1,9 \\
G2\_homem\_mulher & — & nnls & 88,1 & 91,9 & 96,2 & +8,1 & 76,4 & 74,5 & 69,8 & $-$6,7 \\
\midrule
G2\_homem\_mulher\_classesAB & — & perceptron & 70,5 & 72,4 & 71,4 & +1,0 & 22,9 & 24,8 & 21,9 & $-$1,0 \\
G2\_homem\_mulher\_classesAB & — & nnls & 72,4 & 77,1 & 75,7 & +3,3 & 31,4 & 36,2 & 31,0 & $-$0,5 \\
\bottomrule
\multicolumn{11}{p{0.95\linewidth}}{\footnotesize Síntese cruzada: o ganho de R$^2$ para R$^4$ indica se o LLM adota um critério não linear. A\_linear = Problema A (rótulos zero-shot); F = meia-lua no pipeline externo (105 pontos de treino); G1, G2 = pipeline externo peso $\times$ altura (Fase A zero-shot). Fidelidade (vs LLM) e acurácia (vs rótulo real) da métrica, contexto. Execução anterior a 10/10/2026 (sem coluna feature\_names): cada linha é a média sobre as variantes de nomes de atributo e de classe do subproblema. Médias das sementes.} \\
\end{tabular}
\end{table}
